\begin{frame}[shrink=4]{Multi-Hop 802.11 Ad Hoc Networks}

\KeyPoint{In ad hoc mode, stations communicate without an access point and relay packets over multiple wireless hops.}

\vspace{0.18cm}

\begin{columns}[T,onlytextwidth]
\column{0.55\textwidth}
\begin{center}
\begin{tikzpicture}[
    station/.style={circle,draw=ResearchColor,fill=CardBg,minimum size=0.82cm,font=\bfseries},
    link/.style={draw=ResearchColor,line width=1pt},
    route/.style={-{Latex[length=2.2mm]},draw=WarningColor,line width=1.5pt}
]
    \node[station] (a) at (0,0) {A};
    \node[station] (b) at (1.65,0.65) {B};
    \node[station] (c) at (3.30,0) {C};
    \node[station] (d) at (4.95,0.65) {D};

    \draw[link] (a) -- (b);
    \draw[link] (b) -- (c);
    \draw[link] (c) -- (d);
    \draw[route] ([yshift=-0.18cm]a.east) -- ([yshift=-0.18cm]b.west);
    \draw[route] ([yshift=-0.18cm]b.east) -- ([yshift=-0.18cm]c.west);
    \draw[route] ([yshift=-0.18cm]c.east) -- ([yshift=-0.18cm]d.west);

    \node[font=\fontsize{8}{9}\selectfont,text=SecondaryColor,align=center] at (2.48,-0.95)
    {One end-to-end packet requires\\three channel-access events};
\end{tikzpicture}
\end{center}

\column{0.41\textwidth}
\fontsize{9}{10.8}\selectfont
\textbf{Key characteristics}
\begin{itemize}
    \item No centralized scheduler or fixed infrastructure.
    \item Nodes act as both end hosts and packet forwarders.
    \item Every hop consumes airtime in a shared broadcast medium.
    \item Neighboring hops cannot always transmit simultaneously.
\end{itemize}
\end{columns}

\vspace{0.12cm}
\SmallNote{\textbf{Implication:} End-to-end throughput depends on route length and local contention around every forwarding node.}

\end{frame}
